\documentclass[12pt,letterpaper, onecolumn]{exam}
\usepackage{amsmath}
\usepackage{amssymb}
\usepackage[lmargin=71pt, tmargin=1.2in]{geometry}  %For centering solution box
\usepackage[brazil]{babel}
\usepackage{natbib}   % very common in economics
\usepackage{enumitem}
\usepackage[hidelinks]{hyperref}
\urlstyle{same}
\usepackage{booktabs}

% \chead{\hline} % Un-comment to draw line below header
\thispagestyle{empty}   %For removing header/footer from page 1

\begin{document}

\begingroup  
    \centering
    \LARGE Exercícios no R - Econometria I - 2026\\[0.5em]
    %\large \today\\[0.5em]
    \large Prof: André Luiz Pereira Mancha \par
    \large Monitor: Tuffy Licciardi Issa  \par
   \small \href{mailto:tuffy@usp.br}{tuffy@usp.br} \quad
\href{https://github.com/Tuffyli}{github.com/tuffyli} \par
    %\large Roll Number\par
    %\large Class/Batch/Section\par
\endgroup
\rule{\textwidth}{0.4pt}
\pointsdroppedatright   %Self-explanatory
\printanswers
\renewcommand{\solutiontitle}{\noindent\textbf{Ans:}\enspace}   %Replace "Ans:" with starting keyword in solution box


\begin{questions}



    \question (2.C10 \citep{wooldridge2016}) O conjunto de dados do arquivo \textit{CATHOLIC} inclui informações de pontuações de testes de mais de 7.000 estudantes dos Estados Unidos que cursaram a oitava série em 1988. As variáveis \textit{math12} e \textit{read12} são notas padronizadas de matemática e leitura, respectivamente.
            \begin{enumerate}[label = (\alph*)]
            \item Quantos estudantes existem na amostra? Encontre as médias e os desvios-padrão de \textit{mate12} e \textit{leitu12}.
            \item Compute a regressão simples de \textit{mate12} sobre \textit{leitu12} para obter o intercepto MQO e as estimativas da inclinação. Reporte os resultados na forma: \begin{equation*}
                \widehat{math12} = \hat\beta_0 + \hat\beta_1read12
            \end{equation*} \begin{equation*}
                n = \ ?, \ \ \ \ \  R^2 = \     ?
            \end{equation*} em que você vai preencher os valores de $\hat\beta_0$ e $\hat\beta_1$, além de substituir os pontos de interrogação.
            \item O intercepto registrado em (b) tem uma interpretação significativa? Explique.
            \item Você está surpreso pelo $\hat\beta_1$ encontrado? E quanto ao $R^2$?
            \item Suponha que você apresente suas descobertas ao superintendente distrital de educação e ele diga "\textit{Suas descobertas mostram que, para aumentar as notas de matemática, precisamos somente melhorar as notas de leitura; portanto, devemos contratar mais professores de leitura}." Como você responderia a esse comentário?(Dica: Se você calculasse a regressão de \textit{read12} sobre \textit{math12}, em vez do contrário, o que esperaria descobrir?)
            \end{enumerate}

    \question (2.C11 \citep{wooldridge2016}) Utilize os dados do GPA1 para responder estas perguntas. É uma amostra de estudantes universitários da Michigan State University a partir de meados dos anos 1990, e inclui o atual GPA universitário, \textit{colGPA}, e uma variável binária que indica se o estudante possuía um computador pessoal (PC).
    \begin{enumerate}[label = (\alph*)]
        \item Quantos alunos são incluídos na amostra? Encontre a média e o maior GPAs da faculdade.
        \item Quantos estudantes possuíam o seu próprio PC?
        \item Estime a equação da regressão simples \begin{equation*}
            colGPA = \beta_0 + \beta_1 PC + u
        \end{equation*} e informe as suas estimativas para $\beta_0$ e $\beta_1$. Interprete essas estimativas, incluindo uma discussão das magnitudes.
        \item O que é R-quadrado da regressão? O que pensa da sua magnitude?
        \item O seu achado no item (c) implica que possuir um PC tem um efeito causal na \textit{colGPA}? Explique.
    \end{enumerate}
    
\end{questions}
\bibliographystyle{apalike}   % or plainnat, econ, etc.
\bibliography{refs}
\end{document}